\begin{abstractpage}

\begin{abstract}{romanian}
Lucrarea de față prezintă proiectarea, implementarea și evaluarea sistemului \textbf{TriSense}, un sistem cyber-fizic adaptiv destinat roboticii asistive în intervențiile pentru copii neurodivergenți. TriSense integrează trei niveluri de execuție interconectate: un robot construit pe platforma LEGO\textsuperscript{®} SPIKE\texttrademark{} Prime (firmware Pybricks), o placă LMS-ESP32 cu MicroPython responsabilă de comunicație și audio, și o aplicație de tip „creier” care rulează pe PC, scrisă în Python. Comunicarea între componente folosește protocolul MQTT, un canal serial LPF2 între hub și ESP32, precum și fluxuri audio PCM prin TCP. Componenta de inteligență artificială se bazează pe modelele Google Gemini pentru dialog, transcrierea vorbirii și analiza vizuală a construcțiilor LEGO realizate de copil, capturate cu un modul ESP32-CAM. Sistemul oferă un set de activități terapeutice interactive — recunoașterea emoțiilor, imitarea tiparelor de mișcare, jocul de construcție „Build the Model” inspirat de Turnul din Hanoi, jocuri de funcții executive (Stop \& Go, estimarea timpului), o poveste co-construită și exerciții de respirație pentru autoreglare emoțională — împreună cu înregistrarea automată de metrici de sesiune utile terapeuților. Rezultatele obținute arată că o arhitectură distribuită, construită din componente accesibile, poate susține interacțiuni multimodale (voce, mișcare, vedere artificială) adaptate nevoilor copiilor neurodivergenți.
\end{abstract}

\begin{abstract}{english}
This thesis presents the design, implementation and evaluation of \textbf{TriSense}, an adaptive cyber-physical system for assistive robotics in interventions for neurodivergent children. TriSense integrates three interconnected execution levels: a robot built on the LEGO\textsuperscript{®} SPIKE\texttrademark{} Prime platform (Pybricks firmware), an LMS-ESP32 board running MicroPython that handles communication and audio, and a Python ``brain'' application running on a PC. The components communicate through the MQTT protocol, an LPF2 serial channel between the hub and the ESP32, and PCM audio streams over TCP. The artificial-intelligence component relies on Google Gemini models for dialogue, speech transcription and visual analysis of the LEGO constructions built by the child, captured with an ESP32-CAM module. The system offers a set of interactive therapeutic activities --- emotion recognition, movement-pattern imitation, the ``Build the Model'' construction game inspired by the Tower of Hanoi, executive-function games (Stop \& Go, time estimation), a co-constructed story, and breathing exercises for emotional self-regulation --- together with the automatic logging of session metrics useful to therapists. The results show that a distributed architecture, built from affordable components, can support multimodal interactions (voice, movement, computer vision) tailored to the needs of neurodivergent children.
\end{abstract}

\end{abstractpage}
